\documentclass{elegantpaper}
\usepackage[utf8]{inputenc}
\usepackage[T1]{fontenc}
\usepackage{geometry}
\usepackage{graphicx}
\usepackage{float}
\usepackage{hyperref}
\usepackage[english,indonesian]{babel}

\title{Optimasi Pemrosesan Template LaTeX Menggunakan Pendekatan Berbasis AI}
\author{Academic Support Agent}
\date{\today}

\begin{document}
\maketitle

\begin{abstract}

\end{abstract}

\chapter{PENDAHULUAN}


\chapter{TINJAUAN PUSTAKA}


% Original content (commented out):
% %!TEX program = xelatex
% 本模板为上海交通大学研究生学位论文开题报告模板，基于 ElegantPaper 修改而来。
% CopyRight: School of Mathematical Sciences, Shanghai Jiao Tong Univeristy.
% 完整编译: xelatex -> biber/bibtex -> xelatex -> xelatex
\documentclass[lang=cn, 11pt, a4paper, AutoFakeBold]{elegantpaper}
\usepackage{caption}
\usepackage{subcaption}
\usepackage{setspace}
\usepackage{fontspec}
\usepackage{setspace}
\usepackage{graphicx}
\usepackage{fancyhdr}
\usepackage{pdfpages}
\usepackage{setspace}
\usepackage{booktabs}
\usepackage{multirow}
\usepackage{wasysym}


\title{PhD Thesis Proposal}
\author{\kaishu{张郁林} \\ \kaishu{数学科学学院}}
\institute{上海交通大学}

\version{0.10}
\date{\zhtoday}


% 本文档命令
\usepackage{array}
\newcommand{\ccr}[1]{\makecell{{\color{#1}\rule{1cm}{1cm}}}}
\usepackage[thicklines]{cancel}

\numberwithin{equation}{section}

\begin{document}

\input{cover.tex} %请到 cover.tex 里修改封面个人信息

\includepdf{instruction-1.pdf} %插入填报说明
\includepdf{instruction-2.pdf}

%\maketitle

%\begin{abstract}
%本文开题报告。
%\keywords{thesis proposal}
%\end{abstract}

\newpage
\begin{center}
\begin{tabular}{|l|l|}
\hline
   \large{论文题目}  & \multirow{2}{*}{\makebox[100mm][l]{\large\bfseries\fangsong 理科楼 6 号楼工位安排与管理}}\\

   \large{Proposed Title}  & \\
\hline
   \multirow{4}{*}{\large{研究课题来源}}  &   {\large\fangsong 请在合适选项前画 \checkmark\,\, Please select proper options by “\checkmark”}\\
                                 &   {\large\fangsong \Square\,\, 国家自然科学基金课题 NSFC Research Grants}\\
                                 &   {\large\fangsong \Square\,\, 国家社会科学基金 National Social Science Fund of China}\\
                                 &  {\large\fangsong \Square\,\, 国家重大科研专项 National Key Research Projects}\\
    \large{Source of Research}   &   {\large\fangsong \Square\,\, 其它纵向科研课题 Other Governmental Research Grants}\\
    \large{Project}             &   {\large\fangsong \Square\,\, 企业横向课题 R$\&$D Projects from Industry}\\
    \multirow{2}{*}{}           &   {\large\fangsong \Square\,\, 自拟课题 Self-proposed Project}\\
                                 &  {\large\fangsong \Square\,\, 其他 Other \underline{\makebox[70mm][l]{\large\fangsong }}}\\
\hline 
\end{tabular}
\end{center}

\vspace{6pt}

\section{请综述课题国内外研究进展、现状、挑战与意义，可分节描述。博士生不少于10,000汉字，硕士生不少于5,000汉字。请在文中标注参考文献\quad\normalsize{Please review the frontier, current status, challenges and significance of the research topic. The citations should be marked in the context and listed in order at the end of this section. No less than 8,000 words for doctoral students and 4,000 words for master students if written in English.}}

{\kaishu %楷体从此开始
\subsection{序言\quad\normalsize{Introduction}}

此处序言

\begin{figure}[htbp]
	\centering
	\includegraphics[scale=0.5]{dewetting-illustration.jpg}
	\caption{固体脱湿现象示意图，其中基底 (绿色) 与薄膜 (蓝色) 均处于固态. 该图取自\cite{thompson2012solid}.}\label{dewetting-illustration}
\end{figure}


%--------------------SECTION-----------------------------%

\subsection{物理背景与界面/表面动力学\quad\normalsize{Physics background and interface/surface dynamics}}

\begin{figure}[htbp]
	\centering
	\begin{subfigure}{0.45\linewidth}
		\centering
		\includegraphics[width=0.9\linewidth]{image/single-crystal-dewetting.jpg}
		\caption{不同尺寸的单晶薄膜的脱湿现象.}\label{single-crystal-dewetting}
	\end{subfigure}
	\centering
	\begin{subfigure}{0.45\linewidth}
		\centering
		\includegraphics[width=0.9\linewidth]{image/aniso-crystal-dewetting.jpg}
		\caption{三种不同平面定向的正方形的脱湿现象.}\label{aniso-crystal-dewettig}
	\end{subfigure}
 \caption{图像均由 Ye 与 Thompson 复现制作 \cite{bao2012mathematical}.}
 \end{figure}

\subsubsection{界面能量与表面扩散\quad\small{Interfacial energy and surface diffusion}}

\begin{theorem}
type theorem here.
\begin{equation}
     W=\frac{1}{c_F}\int_{\Omega}\varepsilon A(\nabla\phi)+\varepsilon^{-1}F(\phi)\,\mathrm{d}x+\sigma\int_{\Gamma_w}G(\phi)\,\mathrm{d}s
 \end{equation}
\end{theorem}

\begin{lemma}
    type lemma here.
    \begin{subequations}\label{phase-field-model-iso}
    \begin{align}
        &\frac{\partial\phi}{\partial t}=\nabla\cdot(M\nabla\mu)\\
        &\mu=F^\prime(\phi)-\varepsilon^2\Delta\phi
    \end{align}
\end{subequations}
\end{lemma}

\begin{remark}
    remark.
\end{remark}

} %楷体结束

\section{课题研究目标、主要研究内容和拟解决的关键问题\quad\normalsize{ Research objectives, main contents and key issues to be solved.}}

{\kaishu %楷体从此开始

t.b.a

\section{拟采取的研究方法、研究方案及其可行性分析\quad\normalsize{Research methods and research scheme to be adopted and feasibility analysis.}}

t.b.a

\section{课题的创新点\quad\normalsize{Novelties of the proposed topic.}}

t.b.a

\section{计划进度、预期成果\quad\normalsize{Research schedule, and expected outcomes.}}

t.b.a

\section{与本课题有关的工作积累、已有的研究工作成绩\quad\normalsize{Prior experience and accomplished achievements related to the proposed topic.}}

t.b.a

} %楷体结束

\section*{本人承诺：开题报告中的内容真实无误，若有不实，愿承担相应的责任和后果\quad\normalsize{I hereby declare and confirm that the details provided in this Form are valid and accurate. If anything untruthful found, I will bear the corresponding liabilities and consequences.}}

\textbf{学生签字/Signature of Student:}\includegraphics[height=1.5cm]{image/Auto.png}\quad\quad\textbf{日期/Date:}\quad\zhtoday %插入你的电子签名

\nocite{*}
\printbibliography[heading=bibintoc, title=\ebibname]

\appendix
%\appendixpage
\addappheadtotoc

\end{document}


\end{document}